\subsection{Pipeline de prétraitement}

Le prétraitement est une étape cruciale qui conditionne la qualité des représentations.
Nous appliquons un pipeline en 6 étapes séquentielles spécialement adapté aux tweets :

\begin{enumerate}[leftmargin=*,label=\textbf{Étape \arabic*.}]
  \item \textbf{Expansion des contractions} : \textit{I'm} $\rightarrow$ \textit{I am}, \textit{don't} $\rightarrow$ \textit{do not}
  \item \textbf{Gestion des emojis} : conversion en texte descriptif (\texttt{😊} $\rightarrow$ \texttt{smiling face})
  \item \textbf{Nettoyage général} :
    \begin{itemize}
      \item URLs $\rightarrow$ token \texttt{URL}
      \item @mentions $\rightarrow$ token \texttt{USER}
      \item \#hashtags $\rightarrow$ extraction du mot (\texttt{\#amazing} $\rightarrow$ \texttt{amazing})
      \item Chiffres $\rightarrow$ token \texttt{NUM}
      \item Suppression de la ponctuation et mise en minuscules
    \end{itemize}
  \item \textbf{Tokenisation} : \texttt{TweetTokenizer} de NLTK (adapté aux tweets)
  \item \textbf{Suppression des stopwords} : liste NLTK anglaise, tokens $\leq$ 2 caractères supprimés
  \item \textbf{Lemmatisation} : \texttt{WordNetLemmatizer} de NLTK
\end{enumerate}

\subsection{Exemple avant/après}

\begin{table}[H]
\centering
\caption{Transformation d'un tweet par le pipeline complet}
\begin{tabularx}{\textwidth}{lX}
\toprule
\textbf{Étape} & \textbf{Texte} \\
\midrule
Original        & \texttt{I'm SO happy today!! 😊 Check this http://t.co/xyz \#amazing @john} \\
\midrule
Contractions    & \texttt{I am SO happy today!! 😊 Check this http://t.co/xyz \#amazing @john} \\
Emojis          & \texttt{I am SO happy today!! smiling face Check this http://t.co/xyz \#amazing @john} \\
Nettoyage       & \texttt{am so happy today smiling face check this URL amazing USER} \\
Tokenisation    & \texttt{[am, so, happy, today, smiling, face, check, this, url, amazing, user]} \\
Stopwords       & \texttt{[happy, today, smiling, face, check, url, amazing]} \\
\textbf{Final}  & \textbf{\texttt{happy today smile face check url amaze}} \\
\bottomrule
\end{tabularx}
\end{table}

\subsection{Résultats du prétraitement}

\begin{figure}[H]
\centering
\includegraphics[width=\textwidth]{figures/02_preprocessing.png}
\caption{Résultats du prétraitement : top mots par classe, distribution des tokens avant/après,
et impact de chaque étape sur le nombre moyen de tokens.}
\label{fig:preprocessing}
\end{figure}

Le prétraitement réduit en moyenne le nombre de tokens de \textbf{40\%}, éliminant le bruit
et les termes non informatifs. Les mots les plus fréquents confirment la cohérence des données :
\textit{sad, hate, miss} dominent les tweets négatifs, tandis que \textit{happy, love, thank}
caractérisent les tweets positifs.
